\section{The field, the graph, and the local group}
\label{op:finite-section}

We first construct the finite data. This fixes every coefficient field
and denominator before the modular argument begins.

\begin{lemma}\label{op:field}
For $m,p$ in Theorem~\ref{op:main}, $p>m$ is an odd prime and $-1$
is a square in $\F_p$. If $\eta$ is a nonsquare in $\F_p^*$, then
\[
 K_0=\F_p[\xi]/(\xi^2-\eta),\qquad
 K=K_0[\nu]/(\nu^2-\xi)
\]
are quadratic field extensions, so $|K|=p^4$.
\end{lemma}
\begin{proof}
A least divisor greater than one of an integer is prime: any nontrivial
proper divisor of a composite least divisor would be smaller. If $p\le m$,
then $p$ divides both $m!$ and $(m!)^2+1$, which is impossible. Hence
$p>m>600$, and the defining divisibility gives $(m!)^2=-1$ in $\F_p$.
There are nonsquares because each nonzero square has exactly two square
roots. The first quotient is therefore a quadratic field. The conjugate
of $\xi$ is $-\xi$, so its norm is $-\eta$. Since $-1$ is square,
$-\eta$ is nonsquare. A square in $K_0$ would have square norm, proving
that $\xi$ is nonsquare and that the second quotient is a field.
\end{proof}

Choose the least nonsquare using representatives $0,\ldots,p-1$; order
$K$ by its coefficients in $(1,\xi,\nu,\xi\nu)$. Put $M=p^4-1$ and
choose the least element $\zeta\in K^*$ of order $M$. This element
exists: if $b$ is the least common multiple of all orders in $K^*$,
then $b\mid M$, while the root bound for $X^b-1$ gives $M\le b$.
For each prime-power factor of $b$, an appropriate power of one element
has that exact order. Multiplying the resulting commuting elements of
coprime orders gives order $b=M$. In particular, $M$ and all its powers
are invertible in $K$.

Use natural orders on integer indices, increasing-list lexicographic
orders on subsets, and lexicographic orders on tuples and finite function
tables. Order characters of $(\mathbb Z/M\mathbb Z)^D$ by their exponent
lists on the ordered coordinate generators. These conventions specify
the finite choices below.

Let $I_1,\ldots,I_m$ be the $600$-subsets of $[1200]$ in order. Join
distinct $i,j$ in $\Gamma$ exactly when $|I_i\cap I_j|<200$. Put
\[
 \mathcal T=\{T\subseteq[1200]:|T|\le400\},\qquad
 t=\sum_{j=0}^{400}\binom{1200}{j}.
\]

\begin{proposition}\label{op:graph}
The graph $\Gamma$ is triangle-free and has a specified fitting
factorization $F_{ij}=\sum_{T\in\mathcal T}\alpha_i(T)\beta_j(T)$
over $\F_p$ with $t/m<1/100$.
\end{proposition}
\begin{proof}
Three pairwise adjacent sets would have union of size at least
$3\cdot600-3\cdot199=1203$, exceeding $1200$. To construct the fit, let
$\alpha_i(T)=1$ when $T\subseteq I_i$, and zero otherwise. Define
$\beta_j(T)$ as the coefficient of $\prod_{x\in T}X_x$ obtained by
expanding
\[
 Q_j(X)=\prod_{k=200}^{599}
       \frac{\sum_{x\in I_j}X_x-k}{600-k}
\]
and replacing every positive exponent in each monomial by one. All
denominators are in $\{1,\ldots,400\}$, hence invertible. This replacement
preserves evaluations on zero-one vectors and does not increase degree.
Thus
\[
 F_{ij}=\prod_{k=200}^{599}\frac{|I_i\cap I_j|-k}{600-k}.
\]
It equals one for $i=j$. For distinct nonneighbors the intersection
size lies between $200$ and $599$, so one factor is zero. This is the
threshold-Kneser polynomial specialization of
\cite[Proposition~13]{Haviv2018Theta}
\cite[Theorem~3.2]{GolovnevHaviv2022}.

The binomial coefficients increase up to index $400$. Consequently
\[
 \frac{t}{m}
 \le401\prod_{u=401}^{600}\frac{u}{1201-u}
 \le401(5/7)^{100}<\frac1{100}.
\]
In the second inequality the first hundred factors are at most
$500/701<5/7$ and the rest are at most one. The last inequality follows
from $(5/7)^3<1/2$ and $40100<2^{33}$.
\end{proof}

For the local algebra define the additive group $P=(\F_p)^3$, the
function group $A_0=(\mathbb Z/M\mathbb Z)^P$, and
\[
 H=A_0\rtimes P,\qquad
 (a,s)(a',s')=(a+s\cdot a',s+s'),\qquad
 (s\cdot a)(x)=a(x-s).
\]
Its order is $w=M^{p^3}p^3$. Let $c\in A_0$ have value one at zero and
zero elsewhere, and put
\[
 u(s)=(0,s),\quad U=u(P),\qquad
 v(s)=cu(s)c^{-1},\quad V=cUc^{-1}.
\]
The coordinate formula $v(s)=(c-s\cdot c,s)$ gives $U\cap V=\{1\}$,
because $c-s\cdot c\ne0$ when $s\ne0$. We next construct an idempotent
with incoming annihilation by $U$ and outgoing annihilation by $V$.
